\section{Conclusions}
The calculations permitted reconstruction of equilibrium populations conditional on rest masses, HRG inventory, temperature, and cosmological excess. At the 50 MeV literature reference, $2.7261\times10^{30}$ baryons and $2.7172\times10^{30}$ antibaryons per cm$^3$ were obtained, with an excess of $8.8765\times10^{27}$ cm$^{-3}$. A total of 306.1147 antibaryons coexisted per unit of net excess; nucleons and their antiparticles contributed 95.2176\% of the total.

The criterion $n_{\bar B}=n_B-n_{\bar B}$ was reached at 38.203739 MeV with PDG2025 and rounded to the published reference of 38.2 MeV. Variation of $\eta$ alone shifted the crossing to 38.192814--38.214598 MeV. The comparison supports consistency of equilibrium models with related conditions without determining a total physical uncertainty.

Present abundance alone did not identify a unique initial population or an annihilation history. The approximation $\mu_Q=\mu_S=0$ left nonzero net charge and strangeness. These limits define the conclusion: a quantitative and reproducible thermal reconstruction was achieved, and a complete cosmological reconstruction requires additional chemical balances and kinetic evolution.

\section*{Data and code availability}
The accompanying package includes manuscript sources, figures for each language, fixed catalogs and license, 18 aggregate states for three lists, 1329 particle-species records for PDG2025, the crossing curve, sensitivity to $\eta$, 24 numerical checks, and a Windows execution guide. Cosmological conversions and the previous analysis are retained as identified background material. The calculation was executed in Python using an independent implementation; Thermal-FIST data were used without executing its C++ program.

\section*{Assistance in preparation}
Artificial intelligence assistance was used to prepare text, code, figures, and layout. The reported calculations and checks were executed and their outputs retained. The PowerShell launcher was reviewed statically; it was not executed on Windows. Scientific interpretation and the publication decision are the author's responsibility.
